\chapter{Annexe B — Commandes utiles (.NET CLI)}
\label{ann:cmd}

\section{Restaurer et compiler}

\begin{verbatim}
cd "chemin\vers\GestionPortefeuille2"
dotnet restore
dotnet build GestionPortefeuille.sln
\end{verbatim}

\section{Exécuter l'application Web}

\begin{verbatim}
cd GestionPortefeuille.Web
dotnet run
\end{verbatim}

Mode développement avec rechargement à chaud :

\begin{verbatim}
dotnet watch run
\end{verbatim}

\section{Migrations Entity Framework Core}

Depuis la racine de la solution (adapter les chemins si besoin) :

\begin{verbatim}
dotnet ef migrations add NomMigration ^
  --project GestionPortefeuille.Infrastructure ^
  --startup-project GestionPortefeuille.Web

dotnet ef database update ^
  --project GestionPortefeuille.Infrastructure ^
  --startup-project GestionPortefeuille.Web
\end{verbatim}

\textit{Remarque :} sous PowerShell, remplacer \verb|^| par le caractère d'échappement de continuation de ligne adéquat ou écrire la commande sur une seule ligne.

\section{Problème de fichiers verrouillés (Windows)}

Si le build échoue avec un message indiquant qu'un \texttt{.dll} ou \texttt{.pdb} est utilisé par un autre processus :

\begin{enumerate}[leftmargin=*]
  \item Arrêter l'application (\texttt{Ctrl+C} dans le terminal où \texttt{dotnet watch} ou \texttt{dotnet run} tourne).
  \item Si nécessaire : \texttt{dotnet build-server shutdown}, puis relancer \texttt{dotnet build}.
\end{enumerate}

\section{Compiler le PDF du rapport}

\begin{verbatim}
cd Rapport
pdflatex -interaction=nonstopmode rapport.tex
pdflatex -interaction=nonstopmode rapport.tex
\end{verbatim}

Ou avec \texttt{latexmk} :

\begin{verbatim}
latexmk -pdf -interaction=nonstopmode rapport.tex
\end{verbatim}
